
\chapter{Testing, Debugging, and Scheduling}
\section{Unit Testing and Fuzzing}
\begin{itemize}
    \item Testing transactional data structures for atomicity invariants.
    \item Fuzzing aborts to stress test undo-log correctness.
    \item Invariant-based testing: e.g.\ money conservation in a bank benchmark, counter equality in a shared-counter suite --- the assertion that catches lost updates.
    \item The best invariants are sums: a race-free TM preserves an accumulator; a defect surfaces as a drift in the sum.
    \item Fuzzing `oracle' question: what is the ground truth for a concurrent data structure? (A sequential reference implementation run in a single thread.)
    \item Coverage goal: exercise abort paths (capacity aborts, conflicts, explicit aborts) far more often than real workloads do.
\end{itemize}

\begin{table}[htbp]
\centering
\caption{Testing layers for transactional-memory implementations.}
\label{tab:tm-testing-layers}
\begin{tabular}{p{0.22\linewidth}p{0.26\linewidth}p{0.26\linewidth}p{0.16\linewidth}}
\hline
Layer & Oracle & Bugs exposed & Cost profile \\
\hline
Invariant test & Sum, count, set membership, or sequential reference result & Lost update, torn write-back, non-atomic commit & Cheap; run on every backend \\
Abort fuzzing & Same invariant after injected aborts & Bad undo log, leaked write-set entry, failed rollback & Moderate; many retries \\
Deterministic scheduling & Recorded interleaving plus invariant failure & Rare conflict window, order-sensitive validation bug & Higher; bounded schedules \\
Record-and-replay & Concrete failing execution & Late manifestation of an earlier bad read or write & Debugging cost after failure \\
TLA+ model checking & Abstract safety property & Missing validation case, illegal state transition & State-space limited \\
\hline
\end{tabular}
\end{table}

\section{Worked Example: An Invariant Test for Lost Updates}
\index{invariant test}\index{money conservation}
\label{sec:test-example}

The money-conservation assertion that Chapter~1's running example demands:

\begin{lstlisting}[language=C, caption={An invariant test for the transfer benchmark.}]
TEST(bank, money_is_conserved) {
    int initial = balance_sum();
    spawn(8, worker);            /* 8 threads each doing transfers */
    join_all();
    ASSERT_EQ(initial, balance_sum());   /* sum is invariant */
}
\end{lstlisting}

\begin{itemize}
    \item A race-free TM preserves the accumulator; a defect (lost update, torn write-back) surfaces as a drift in the sum --- Exercise~1 of this chapter.
    \item The same binary runs against every backend in Part~IV with one flag, so one test suite validates the whole family.
    \item Exercise~1 asks for the abort-fuzzing extension; \S\ref{sec:test-example} feeds the oracle question of this section: what is ground truth?
    \item Combine with \S\ref{sec:math-example}: at abort probability
    \[
    p = 0.9
    \]
    eight threads completing ten successful transactions each require the expected number of attempts
    \[
    8 \times 10 \div (1-p) = 800 .
    \]
    The test still passes only if undo logging is correct.
\end{itemize}

\section{Worked Example: Abort Fuzzing for the Bank Invariant}
\index{abort fuzzing}\index{undo log}
\label{sec:abort-fuzz-bank}

The following standalone C++ test uses a single global lock to make the transaction body serial, then injects explicit aborts after partial updates. The point is not performance. The point is coverage of rollback paths that normal executions rarely hit.

\begin{lstlisting}[language=C++, caption={A compilable abort-fuzzing harness for the bank benchmark.}]
#include <cassert>
#include <cstdint>
#include <exception>
#include <iostream>
#include <mutex>
#include <numeric>
#include <random>
#include <thread>
#include <utility>
#include <vector>

struct Abort : std::exception {};

struct Account {
    int balance;
};

static int balance_sum(const std::vector<Account>& bank) {
    return std::accumulate(bank.begin(), bank.end(), 0,
        [](int acc, const Account& a) { return acc + a.balance; });
}

static bool transfer_with_fuzzed_abort(std::vector<Account>& bank,
                                       std::mutex& global_lock,
                                       int from,
                                       int to,
                                       int amount,
                                       std::mt19937& rng,
                                       double abort_probability) {
    std::bernoulli_distribution abort_now(abort_probability);
    std::vector<std::pair<int, int> > undo;  // (index, old_balance)

    auto write_balance = [&](int index, int value) {
        undo.push_back(std::make_pair(index, bank[index].balance));
        bank[index].balance = value;
    };

    global_lock.lock();
    try {
        write_balance(from, bank[from].balance - amount);

        if (abort_now(rng)) {
            throw Abort();
        }

        write_balance(to, bank[to].balance + amount);

        if (abort_now(rng)) {
            throw Abort();
        }

        global_lock.unlock();       // commit
        return true;
    } catch (const Abort&) {
        for (auto it = undo.rbegin(); it != undo.rend(); ++it) {
            bank[it->first].balance = it->second;
        }
        global_lock.unlock();       // abort
        return false;
    }
}

static void worker(std::vector<Account>& bank,
                   std::mutex& global_lock,
                   int id,
                   int successful_transactions) {
    std::mt19937 rng(0xC0FFEEu + static_cast<unsigned>(id));
    std::uniform_int_distribution<int> account(0,
        static_cast<int>(bank.size()) - 1);
    std::uniform_int_distribution<int> amount(1, 7);

    int committed = 0;
    while (committed < successful_transactions) {
        int from = account(rng);
        int to = account(rng);
        if (from == to) {
            continue;
        }

        if (transfer_with_fuzzed_abort(bank, global_lock, from, to,
                                       amount(rng), rng, 0.90)) {
            ++committed;
        }
    }
}

int main() {
    std::vector<Account> bank(64, Account{1000});
    std::mutex global_lock;

    const int initial = balance_sum(bank);

    std::vector<std::thread> threads;
    for (int i = 0; i != 8; ++i) {
        threads.emplace_back(worker, std::ref(bank), std::ref(global_lock),
                             i, 10);
    }

    for (std::thread& t : threads) {
        t.join();
    }

    assert(initial == balance_sum(bank));
    std::cout << "money conserved\n";
    return 0;
}
\end{lstlisting}

\begin{itemize}
    \item The injected abort after the debit catches a missing undo entry for the source account.
    \item The injected abort after the credit catches a rollback loop that restores only the first write.
    \item A global lock is acceptable in the harness: the test is aimed at rollback correctness, not scalability.
    \item The same shape appears in production TM tests: run the normal workload, raise the abort rate, assert the invariant after all threads join.
\end{itemize}

\section{Record-and-Replay and Deterministic Scheduling}
\begin{itemize}
    \item Using \texttt{rr} (Mozilla) for reverse debugging of TM executions.
    \item Deterministic schedulers (CHESS) to systematically explore interleavings.
    \item Why reverse debugging matters: a transaction's bug manifests at commit, long after the offending read.
    \item Deterministic replay as a TM-native debugging tool: replay the exact interleaving that produced a lost update.
    \item The bound on preemptions: exploring all bounded-preemption schedules for small bounds is tractable; full interleaving space is not.
    \item Complementarity with model checking: TLA+ explores the abstract state space; \texttt{rr}\slash{}CHESS explore the concrete one.
\end{itemize}

\begin{figure}[htbp]
\centering
\begin{tikzpicture}[
    node distance=1.6cm and 1.9cm,
    box/.style={draw, rounded corners, align=center, minimum width=2.7cm, minimum height=0.9cm},
    arrow/.style={-{Latex[length=2mm]}, thick}
]
\node[box] (seed) {fuzz seed\\or schedule};
\node[box, right=of seed] (run) {run TM\\backend};
\node[box, right=of run] (check) {check\\invariants};
\node[box, below=of check] (record) {record failing\\interleaving};
\node[box, left=of record] (replay) {deterministic\\replay};
\node[box, left=of replay] (fix) {watchpoint\\and patch};

\draw[arrow] (seed) -- (run);
\draw[arrow] (run) -- (check);
\draw[arrow] (check) -- node[right]{fail} (record);
\draw[arrow] (record) -- (replay);
\draw[arrow] (replay) -- (fix);
\draw[arrow] (fix) |- ($(run.south)+(0,-0.35)$) -- (run.south);
\draw[arrow] (check.north) to[out=90,in=90,looseness=1.1] node[above]{pass: next seed} (seed.north);
\end{tikzpicture}
\caption{A TM debugging loop: random or systematic exploration finds a failing interleaving; replay makes the failure local enough to debug.}
\label{fig:tm-debugging-loop}
\end{figure}

\section{Worked Example: Replay of a Lost-Update}
\index{deterministic scheduling}\index{replay}
\label{sec:replay-lost-update}

The next program is deliberately not a correct TM implementation. It is a minimal deterministic scheduler that reproduces a money-conservation failure by forcing a particular interleaving. Each account is atomic to avoid undefined behavior in the C++ memory model; the bug is logical non-atomicity, not a data race.

\begin{lstlisting}[language=C++, caption={A deterministic replay schedule that exposes a non-atomic transfer.}]
#include <atomic>
#include <cassert>
#include <condition_variable>
#include <iostream>
#include <mutex>
#include <thread>
#include <utility>
#include <vector>

class Replay {
public:
    explicit Replay(std::vector<std::pair<int, int> > script)
        : script_(std::move(script)), pos_(0) {}

    void checkpoint(int tid, int step) {
        std::unique_lock<std::mutex> lk(mu_);
        cv_.wait(lk, [&] {
            return pos_ < script_.size()
                && script_[pos_] == std::make_pair(tid, step);
        });
        ++pos_;
        cv_.notify_all();
    }

private:
    std::vector<std::pair<int, int> > script_;
    std::size_t pos_;
    std::mutex mu_;
    std::condition_variable cv_;
};

struct Bank {
    std::atomic<int> a{100};
    std::atomic<int> b{100};
};

static void nonatomic_transfer(Bank& bank, Replay& replay,
                               int tid, int amount) {
    replay.checkpoint(tid, 0);
    int old_a = bank.a.load();

    replay.checkpoint(tid, 1);
    int old_b = bank.b.load();

    replay.checkpoint(tid, 2);
    bank.a.store(old_a - amount);

    replay.checkpoint(tid, 3);
    bank.b.store(old_b + amount);
}

int main() {
    Bank bank;

    /*
       T0 debits A, T1 observes that debit, T1 credits B,
       then T0 overwrites B with an older value.
    */
    Replay replay({
        {0, 0}, {0, 1}, {0, 2},
        {1, 0}, {1, 1}, {1, 2}, {1, 3},
        {0, 3}
    });

    std::thread t0(nonatomic_transfer, std::ref(bank),
                   std::ref(replay), 0, 10);
    std::thread t1(nonatomic_transfer, std::ref(bank),
                   std::ref(replay), 1, 20);

    t0.join();
    t1.join();

    const int final_sum = bank.a.load() + bank.b.load();

    assert(bank.a.load() == 70);
    assert(bank.b.load() == 110);
    assert(final_sum == 180);

    std::cout << "replayed invariant failure: sum = "
              << final_sum << "\n";
    return 0;
}
\end{lstlisting}

\begin{itemize}
    \item The schedule is the test case. Once recorded, the failure is no longer a Heisenbug.
    \item Atomics remove undefined behavior from the demonstration; they do not provide transaction atomicity.
    \item A watchpoint on \texttt{bank.b} identifies the overwrite: thread 0 writes a value computed before thread 1's credit.
    \item A real TM backend prevents this execution by locking, validating, versioning, or aborting one transaction.
\end{itemize}

\section{Summary}
\begin{itemize}
    \item Debugging concurrent TM code requires tools that can replay rare interleavings.
    \item Systematic testing can expose Heisenbugs that appear only under specific schedules.
    \item Invariant tests are most useful when they are backend-independent: the same bank, counter, or data-structure oracle should run against SGL, TL2, NOrec, TinySTM variants, SwissTM, persistent backends, distributed backends, and GPU designs.
    \item The companion implementation uses both concrete tests and abstract specifications: the concrete suites check transactional and data-structure behavior, while TLA+ specifications check protocol-level safety.
\end{itemize}

\section{Exercises}
\begin{enumerate}
    \item Design a fuzzing strategy that randomly injects aborts to stress-test a TM's undo-log correctness.
    \item Given a recorded \texttt{rr} trace where a transaction T commits but its writes never appear, hypothesise where to set a watchpoint and how to step backwards to find the bug.
    \item A deterministic scheduler explores interleavings with a bounded number of preemptions. Explain how this bound trades off completeness for tractability.
    \item \emph{[implementation]} Extend Listing~\ref{sec:abort-fuzz-bank}'s harness with account floors: no committed transaction may leave an account below zero. Keep money conservation as a separate invariant, and report which invariant fails first.
    \item \emph{[verification]} Write a two-account TLA+ model of transfer with actions \textsc{Read}, \textsc{Debit}, \textsc{Credit}, \textsc{Commit}, and \textsc{Abort}. Add the money-conservation invariant and then intentionally omit rollback for \textsc{Debit}; confirm that the model checker finds a counterexample.
    \item \emph{[theory]} Bounded-preemption testing often finds bugs with a small bound. Give one reason this empirical fact is plausible for TM implementations, and one class of TM bug that may require a larger bound or a different exploration strategy.
\end{enumerate}
